% Teil II: der durchgerechnete Weg.

\subsection*{Gegeben und gesucht}
Bruchgleichungen (a)--(c) nach $x$ auflösen; Definitionsbereich beachten.

\subsection*{a) $\frac1{x+2}=\frac9x$, mit $x\neq0$, $x\neq-2$} 
Über Kreuz:
\begin{equation}
  x=9(x+2)=9x+18\;\Longrightarrow\;-8x=18\;\Longrightarrow\;x=-\frac94 .
\end{equation}
liegt im Definitionsbereich.

\subsection*{b) $\frac3{2x+1}-\frac2{2-x}=0$, mit $x\neq-\frac12$, $x\neq2$}
\begin{equation}
  \frac3{2x+1}=\frac2{2-x}\;\Longrightarrow\;3(2-x)=2(2x+1)
  \;\Longrightarrow\;6-3x=4x+2\;\Longrightarrow\;x=\frac47 .
\end{equation}
liegt im Definitionsbereich.

\subsection*{c) $\frac{x-2}{3+x}=\frac{2x}{2x-3}$, mit $x\neq-3$, $x\neq\frac32$}
\begin{equation}
  (x-2)(2x-3)=2x(3+x)
\end{equation}
\begin{align}
  2x^2-3x-4x+6&=6x+2x^2, \\
  2x^2-7x+6&=6x+2x^2 \;\Longrightarrow\; -13x+6=0
  \;\Longrightarrow\; x=\frac6{13} .
\end{align}
liegt im Definitionsbereich.

\begin{ergebnis}
\begin{equation*}
  \text{a)}\ x=-\frac94,\qquad
  \text{b)}\ x=\frac47,\qquad
  \text{c)}\ x=\frac6{13}.
\end{equation*}
\end{ergebnis}

\begin{kontrolle}
a) $\frac{1}{-9/4+2}=\frac{1}{-1/4}=-4$, rechts $\frac{9}{-9/4}=-4$;
b) $\frac3{2\cdot4/7+1}=\frac3{15/7}=\frac75$, $\frac2{2-4/7}=\frac2{10/7}=\frac75$;
c) $\frac{6/13-2}{3+6/13}=\frac{-20/13}{45/13}=-\frac49$,
$\frac{2\cdot6/13}{12/13-3}=\frac{12/13}{-27/13}=-\frac49$. Alle stimmen.
\end{kontrolle}

\begin{merksatz}
Bei Bruchgleichungen zuerst Nenner $\neq0$ notieren, über Kreuz
multiplizieren, lösen und die Definitionsbedingungen zum Schluss prüfen.
\end{merksatz}
